\documentclass[11pt,a4paper]{article}

% ---------------------------------------------------------------------------
% Classificacao de topicos - case VDE
% Fonte oficial do relatorio. Compila no Overleaf com pdfLaTeX.
% Um unico arquivo: a figura e inline (TikZ), sem assets externos.
% ---------------------------------------------------------------------------

\usepackage[T1]{fontenc}
\usepackage[utf8]{inputenc}
\usepackage[brazilian]{babel}
\usepackage[margin=2.0cm]{geometry}
\usepackage{booktabs}
\usepackage[table]{xcolor}
\usepackage{microtype}
\usepackage{enumitem}
\usepackage{titlesec}
\usepackage{tcolorbox}
\tcbuselibrary{skins}
\usepackage{tikz}
\usetikzlibrary{arrows.meta,positioning}
\usepackage{fancyhdr}
\usepackage{lastpage}
\usepackage[scaled=1.0]{helvet}
\renewcommand{\familydefault}{\sfdefault}
\usepackage[hidelinks]{hyperref}

\definecolor{azul}{HTML}{1F4E79}
\definecolor{cinza}{HTML}{F2F4F7}
\definecolor{cinzaforte}{HTML}{DCE3EB}

\newcommand{\bd}[1]{\textbf{#1}}
\setlist[itemize]{leftmargin=1.2em,itemsep=1pt,topsep=2pt}
\setlength{\parskip}{4pt}
\setlength{\parindent}{0pt}

% --- cabecalho de secao: caixa alta, navy, com filete ---
\titleformat{\section}{\normalfont\large\bfseries\color{azul}}{}{0pt}{\MakeUppercase}[{\vspace{-0.15em}\color{azul}\titlerule[0.8pt]}]
\titlespacing*{\section}{0pt}{10pt}{4pt}

% --- rodape (pagina X de Y) ---
% "Y" vem do lastpage e so aparece certo apos DUAS compilacoes
% (na 1a passada o LaTeX ainda nao sabe qual e a ultima pagina e escreve "??").
\fancypagestyle{rodape}{%
  \fancyhf{}%
  \renewcommand{\headrulewidth}{0pt}%
  \fancyfoot[L]{\small\color{black!55}Classificação de tópicos --- case VDE}%
  \fancyfoot[R]{\small\color{black!55}Página \thepage\ de \pageref{LastPage}}%
}
\pagestyle{rodape}

% --- cartao de indicador (KPI) ---
\newcommand{\kpi}[2]{%
  \fcolorbox{azul!55}{white}{\parbox[c][13mm][c]{0.285\textwidth}{\centering
    {\Large\bfseries\color{azul}#1}\\[2pt]{\scriptsize\bfseries\color{black!65}#2}}}%
}

\tikzset{
  box/.style={draw=azul,rounded corners,fill=white,align=center,inner sep=3pt,
              minimum height=8mm,text width=19mm,font=\small},
  arr/.style={-{Latex[length=2mm]},thick,draw=azul}
}

\begin{document}
% --- bloco de titulo ---
{\LARGE\bfseries\color{azul}Classificação de tópicos: case VDE\par}
\vspace{2pt}
{\large Do raciocínio ao desenho de produção\par}
\vspace{4pt}
{\small Autor: Léo Vitor\par}
\vspace{4pt}
{\color{azul}\hrule height 0.8pt}
\vspace{7pt}
\thispagestyle{rodape}

% ============================ SUMARIO EXECUTIVO ============================
\begin{tcolorbox}[colback=cinza,colframe=azul!35,
                  borderline west={2.5pt}{0pt}{azul},
                  title={\bfseries SUMÁRIO EXECUTIVO},coltitle=azul,colbacktitle=cinza,
                  boxrule=0.6pt,arc=2pt,left=6pt]
\begin{itemize}
\item \bd{Resultado:} publiquei 222 das 359 questões (61,8\%), mandei 136 para revisão e
      deixei uma fora do escopo.
\item \bd{Qualidade:} o sistema concorda de 0,90 a 0,92 com dois rotuladores de IA, que
      entre si concordam em 0,85. Isso mede concordância com um rotulador, não precisão.
\item Com uma amostra deste tamanho não separo 90\% de 95\%, e por isso não afirmo a meta.
\end{itemize}
\end{tcolorbox}

\vspace{2pt}
\noindent
\kpi{61,8\%}{TAXA DE PUBLICAÇÃO}\hfill
\kpi{0,86}{PREVISÃO FINAL (0,74--0,94)}\hfill
\kpi{0,90--0,92}{CONCORDÂNCIA COM ROTULADOR DE IA}
\vspace{4pt}

Minha previsão de precisão é 0,86, no intervalo de 0,74 a 0,94, que é julgamento e não
bootstrap. A cobertura, de 0,62, é uma contagem do pipeline, não uma estimativa. Baixei de
0,90 porque o número antigo vinha de um único rotulador de IA; com uma segunda avaliação
independente, ele parecia otimista.

% ================================= a) COMO PENSEI =========================
\section*{a) Como eu pensei}

Comecei pelo caminho mais simples: pedir a um modelo que escolhesse o tema de cada questão
entre os temas que uma busca por palavras trazia. As primeiras rodadas mostraram dois
problemas: em várias questões o tema certo nem estava entre os candidatos, e o
modelo escolhia o vizinho mais parecido com toda a confiança; e, ao pôr dois modelos para
votar, eles concordavam justamente nesses erros. A votação, que eu esperava que fosse um
alarme, então confirmava o erro em vez de apontá-lo.

O desenho final tem três partes. A busca por palavras propõe os temas mais parecidos, mais os
pais e os filhos dos mais bem colocados. Dois modelos Gemini (\texttt{gemini-\allowbreak 3.6-\allowbreak flash} e
\texttt{gemini-\allowbreak 3.1-\allowbreak flash-\allowbreak lite}) votam duas vezes cada, num total de quatro votos. Um
verificador, o menor dos dois modelos, enxerga a árvore inteira de temas, dá o quinto voto,
que desempata um 2 a 2, e manda a questão para revisão quando discorda. Só então as regras consolidam: votos iguais publicam; votos em
subtemas do mesmo ramo publicam no tema pai (a régua do case aceita pai e filho); temas
diferentes vão para revisão; e a maioria de \texttt{fora} vira \texttt{fora}.

\begin{tcolorbox}[colback=cinza,colframe=azul!30,boxrule=0.6pt,arc=2pt,
                  title={\small\bfseries\color{azul}ARQUITETURA DO PIPELINE DE CLASSIFICAÇÃO ENTREGUES},
                  coltitle=azul,colbacktitle=cinza,left=4pt,right=4pt,top=6pt,bottom=5pt]
\centering
\begin{tikzpicture}[node distance=6mm]
\node[box] (clean) {Limpeza\\\scriptsize HTML $\to$ texto};
\node[box,right=of clean] (cand) {Candidatos\\\scriptsize TF-IDF + pais/filhos};
\node[box,right=of cand] (juiz) {Juiz\\\scriptsize 2 modelos, 4 votos};
\node[box,below=9mm of juiz] (ver) {Verificador\\\scriptsize 5º voto, árvore inteira};
\node[box,right=of ver] (reg) {Regras\\\scriptsize decide};
\node[box,right=of reg] (out) {publicar / revisar / fora};
\draw[arr] (clean)--(cand);
\draw[arr] (cand)--(juiz);
\draw[arr] (juiz)--(ver);
\draw[arr] (ver)--(reg);
\draw[arr] (reg)--(out);
\end{tikzpicture}
\\[2pt]
{\scriptsize\itshape\color{black!60}O verificador entra antes das regras e vê a árvore inteira.}
\end{tcolorbox}

Dos onze erros que achei entre as publicadas da primeira versão, seis eram de busca, com o
tema certo fora da lista de candidatos, e cinco de decisão, com o tema certo na lista e o
juiz escolhendo outro. O
que me chamou a atenção foi que todos eram votos unânimes: a votação nunca desconfiou. Foi
isso que me levou ao verificador, um quinto voto que enxerga a árvore inteira e manda a
questão para revisão quando discorda.

Antes de rodá-lo, escrevi o critério de sucesso: vetar pelo menos três erros reais e recusar
no máximo um quinto das questões certas. Ele vetou cinco dos onze erros e recusou três das
questões certas, dentro do critério. Dos seis que escaparam, minha leitura é que dois dividem
opinião, sendo um a favor do rótulo de IA e um a favor do modelo, e quatro são ambíguos,
todos registrados em \texttt{docs/adjudicacao\_erros.csv}. É a leitura de um revisor, não de um
especialista.

Cheguei ao desenho final por tentativa e erro, e descartei várias ideias pelo caminho. A primeira foi uma política que
publicava só os casos mais confiáveis: chegou a 95 questões com 100\% de concordância, mas eu
a abandonei porque se apoiava em 17 rótulos da mesma amostra que a calibrou. Buscas
alternativas, como BM25, fusão de rankings e embeddings densos, não superaram a busca simples por
palavras. Subir sempre para o tema pai resolveu um caso só, porque o erro raramente era entre
irmãos. E descartei o cross-encoder sem testar, porque ele só reordena a lista, e o problema é
o tema que não entra nela.

O que funcionou melhor foi atacar a busca offline, sem gastar API. Um dicionário jurídico (que
traduz uma citação de lei, como ``Lei 11.340'', para o nome que aparece no título do tema,
``Lei Maria da Penha''), somado à busca por temas irmãos, subiu a taxa com que a busca acha o
tema certo ou um parente direto de 0,67 para 0,96. Esse teste não foi totalmente cego, porque
ampliei o dicionário olhando os erros que sobraram. Quando levei o dicionário até a decisão
final, numa segunda rotulagem independente, o ganho quase desapareceu: a busca foi de 0,79
para 0,85, mas a concordância das publicadas mudou menos de meio ponto percentual. Os acertos
extras caíam em revisão, porque o juiz achava o tema e a agregação não publicava. Por isso, e
porque mudar a lista de candidatos invalidaria o cache das chamadas, não adotei a mudança na
entrega.

% ================================= b) DIAGNOSTICO =========================
\section*{b) Diagnóstico}

Os problemas que mais atrapalham vêm do catálogo e dos dados, mais do que do modelo.

\bd{Temas gêmeos.} O catálogo tem oito pares de temas com o mesmo título em disciplinas
diferentes, por exemplo Lei Maria da Penha em Legislação Especial e em Penal. A régua do case
aceita só um dos dois como certo, e eu não sei qual; perguntei à organização, que respondeu
que a decisão faz parte do case. Adotei então uma premissa conservadora: publico o gêmeo só
quando a disciplina de origem da questão coincide com a raiz do tema, e mando o resto para
revisão. O motivo é que publicar no tema errado é o erro mais caro (o aluno estudaria o tópico
errado), enquanto revisar só atrasa. A premissa custa 24 questões, cerca de 6,7 pontos
percentuais de cobertura, e deixa de valer se o professor confirmar qual gêmeo é o oficial.

\bd{Dados.} Três coisas confundem tanto a máquina quanto uma pessoa. O assunto de origem
engana: a Q0026, sobre Lei Maria da Penha, veio como ``Ação Penal''. A disciplina de origem é
bem mais confiável e coincide com a do tema escolhido em cerca de 95\% das publicadas
rotuladas, embora parte das questões venha de fora do recorte (Constitucional, Civil,
Tributário, Ambiental) ou sem valor. E há questões que cabem em mais de um tema (Q0016 e
Q0085, por exemplo), nas quais o ``tema certo'' é uma escolha, não um fato.

\bd{Que questões são difíceis.} Vi três perfis: assertivas I a IV com pouca pista no
enunciado; legislação especial, em que o texto da questão não repete o título do tema; e as
que cabem em dois temas. O erro não se concentrou num ramo: dos onze erros da primeira
versão, três eram de legislação especial e o resto se espalhava por Parte Geral, Processo
Penal e ambiguidade. As lixeiras (temas com muitas questões na origem) pareciam um risco, mas
nessas o erro foi menor do que nos demais temas, então a regra extra de confiança que criei
para elas quase não filtra nada.

\bd{Bloqueio por conteúdo sensível.} Uma questão (a Q0227, com cenários de violência sexual)
teve a resposta bloqueada pelo filtro de segurança do modelo. O pipeline trata isso como
abstenção: a questão vai para revisão com o motivo registrado, e a chamada não se repete. Em
produção, isso significa que um filtro sem critério jurídico decide o que cai na fila humana,
e esse tipo de bloqueio pode não ser ajustável. Eu recomendaria contar esses casos e
mantê-los na revisão humana, em vez de tentar afrouxar o filtro.

\bd{Como medi, e o que a medida vale.} Não tenho gabarito, então comparei as decisões com
rótulos feitos por IA: primeiro o DeepSeek, sobre as publicadas, e depois o Opus, para ver se
o resultado se repetia. Não consigo garantir que o DeepSeek não tenha visto saídas do modelo.
A tabela resume o que encontrei.

\begin{center}
\small
\begin{tabular}{lcc}
\toprule
\rowcolor{azul}\color{white}\bfseries Comparação & \color{white}\bfseries Itens & \color{white}\bfseries Resultado \\
\midrule
\rowcolor{cinzaforte}\multicolumn{3}{l}{\emph{Concordância, aceitando pai e filho}} \\
Antes do verificador (261 publicadas) $\times$ DeepSeek & 68 & 0,84 \;(0,73 a 0,91) \\
Entrega (222 publicadas) $\times$ DeepSeek & 60 & 0,90 \;(0,80 a 0,95) \\
Entrega $\times$ Opus (quase fora da amostra) & 88 & 0,92 \;(0,84 a 0,96) \\
DeepSeek $\times$ Opus & 41 & 0,85 \;(0,72 a 0,93) \\
\quad \emph{no tema exato} & 41 & 0,54 \\
\midrule
\rowcolor{cinzaforte}\multicolumn{3}{l}{\emph{Verificador, sobre as 261 publicadas antes dele}} \\
Erros apontados e vetados & 11 & 5 \\
Questões certas recusadas & 57 & 3 \\
\midrule
\rowcolor{cinzaforte}\multicolumn{3}{l}{\emph{Limiar de confiança, para chegar perto de 95\% de precisão}} \\
Melhor precisão com cobertura útil & 15 & 0,93 \;(0,70 a 0,99) \\
Precisão máxima & 6 & 1,00 \;(0,61 a 1,00) \\
\bottomrule
\end{tabular}

\smallskip
\parbox{0.92\textwidth}{\small \emph{Nota:} o Opus rotulou 88 questões, 41 em comum com o
conjunto que usei para desenhar as regras, por isso o chamo de \emph{quase} fora da amostra.}
\end{center}

O achado que mais pesa é a quarta linha: os dois rotuladores concordam em 0,85 quando aceito
pai e filho, mas em 0,54 no tema exato. A precisão que consigo medir depende, então, de quem
rotula, e a de especialistas pode cair em qualquer ponto dessa faixa. Com amostra deste
tamanho, os intervalos de 90\% e de 95\% se sobrepõem; para separá-los com folga eu precisaria
de umas 430 questões por grupo.

Também medi quanto custaria chegar perto de 95\%. Subindo o limiar de confiança, eu só
alcançava precisão alta com cobertura muito baixa, como mostra a última parte da tabela. Com
amostras tão pequenas isso não prova nada, e foi por isso que preferi a entrega com 62\% de
cobertura, dizendo abertamente que não demonstrei a meta.

% ================================= c) PRODUCAO ============================
\section*{c) Produção (50 mil questões por semana)}

Esta seção é uma proposta: não a executei, por falta de orçamento de API e de tempo.

\bd{Encaixe no fluxo.} As questões novas entram pelo importador, e a busca propõe os
candidatos. Um classificador leve, treinado com os rótulos que os professores forem
acumulando, resolve a maioria dos casos e devolve uma probabilidade calibrada. O que sobra vai
para o LLM, e um verificador confere os casos ambíguos. O que ainda divergir vai para a fila
de revisão humana, e os reportes de alunos e as revisões alimentam uma amostra auditada que
mede a qualidade toda semana.

\bd{Do barato para o caro.} Primeiro, a busca ampliada (dicionário jurídico e irmãos), que é
grátis. Depois, o classificador treinado, que só existe quando houver rótulos de professores:
com os rótulos de IA que tenho, um classificador linear acertou o tema entre os dez primeiros
em 44\% dos casos. Por último, o LLM com verificação seletiva e a revisão humana nas
divergências.

\bd{Custo.} Calculei etapa por etapa. O fallback e o teste do dicionário estão medidos em
arquivo; o juiz base eu estimei; e o verificador li no console, porque os tokens não foram
gravados em arquivo.

\begin{center}
\begin{tabular}{lcc}
\toprule
\rowcolor{azul}\color{white}\bfseries Etapa (359 questões) & \color{white}\bfseries Tokens de entrada & \color{white}\bfseries Custo (US\$, preço de tabela) \\
\midrule
Juiz (4 votos) & $\sim$2,1 M & 1,04 \\
Fallback (28 questões $\times$ 4) & $\sim$0,9 M & 0,47 \\
Verificador (246 chamadas) & $\sim$2,14 M & 0,56 \\
\midrule
\rowcolor{cinzaforte}\bfseries Total & $\sim$5,2 M & 2,1 \\
Projeção para 50 mil por semana & --- & US\$ 290 (piso) \\
\bottomrule
\end{tabular}
\end{center}

Os preços são de tabela, e a conta acima é um piso: só tokens de entrada, cerca de 14 milhões
por mil questões. O verificador manda mais tokens por chamada, mas custa menos por questão do
que os quatro votos do juiz, porque usa só o modelo barato.

O gasto real foi maior: a fatura do painel ficou em R\$ 32,98, contando também os
experimentos. Dividido pelas 359 questões, o teto observado é de R\$ 0,09 por questão, contra
um piso de US\$ 0,006 na tabela. Para as 50 mil questões, a projeção vai de US\$ 290, no piso,
a R\$ 4,6 mil, no teto observado.

\bd{Cache de contexto.} O bloco que se repete em toda chamada (instruções mais a árvore, uns
9 mil tokens) pode ser enviado uma vez e reaproveitado. O cálculo em tokens, que não gasta
nada (\texttt{-{}-dry-run}), indica cerca de 72\% de economia na entrada. O teste real foi
pequeno e, como os créditos de API acabaram, não o reexecutei; o comando para reproduzir está
em \texttt{docs/cache\_contexto.md}, e o desconto em reais só se confirma no contrato.

\bd{Limitações.} Juiz e verificador são da mesma família de modelos
(Gemini) e podem errar juntos: a Q0120 passou pelos dois. O limiar de confiança quase não
filtra, porque as confianças ficam perto de 0,95. As regras do verificador e dos gêmeos foram
desenhadas na mesma amostra de 60 questões em que são medidas. E a verificação encarece a
conta, porque envia a árvore inteira em cada chamada. Por fim, tudo o que reporto aqui é
concordância com rotuladores de IA, não precisão contra especialistas.

\bd{Como mediria a qualidade.} Toda semana, com três números: a precisão numa amostra auditada
por professores; a taxa de correção da fila (quantas questões revisadas estavam erradas); e os
reportes de alunos por tema, que mostram erros que a auditoria não pegou. A amostra precisa
ser de centenas de itens, não de dezenas, para que os intervalos separem as metas.

\bd{O que eu faria com mais tempo.} Cinco coisas: um verificador de outra família de modelos
(hoje juiz e verificador são ambos Gemini); rótulos de especialistas para 20 a 30 casos de
desacordo; irmãos na busca só nos casos de baixo risco; o cache de contexto em produção; e a
confirmação da convenção dos gêmeos. Com o professor no circuito, acrescentaria ligações
curadas entre leis e temas e entre gêmeos, uma regra de decisão por custo (quando houver
probabilidade calibrada) e aprendizado ativo: a revisão dele gera rótulos que treinam o
classificador.

% ================================= d) MEMO ================================
\clearpage
\section*{(d) Memo para o professor}
{\small\itshape Este memo é para você, especialista em Direito, não em IA.}

\bd{O que eu fiz}
\begin{itemize}
\item Montei um sistema que decide, para cada uma das 359 questões, se publica, manda para
      revisão ou considera fora do escopo. Publiquei 222 (61,8\%), mandei 136 para revisão e 1
      ficou fora.
\item Comparei com duas avaliações feitas por IA, não por especialistas. O sistema concorda
      com elas de 0,90 a 0,92, mas isso não prova que acerta 95\%.
\end{itemize}

\bd{O que encontrei}
\begin{itemize}
\item Quando erra, o sistema erra com confiança: nos casos que examinei, os dois modelos
      concordaram no tema errado. Em geral isso acontece porque o tema certo nem aparece entre
      as opções que o modelo recebe, e não por falta de conhecimento jurídico. Por isso
      acrescentei uma segunda checagem, que enxerga todos os temas e manda a questão para você
      quando discorda.
\item O catálogo tem armadilhas: oito temas aparecem duas vezes, com o mesmo nome, em
      disciplinas diferentes, e há temas que se sobrepõem (pai e filho dizendo quase a mesma
      coisa). Quando há dúvida entre os nomes iguais, prefiro não publicar. Não vi erro causado
      por isso, mas foram poucos casos.
\item Uma busca melhor (dicionário jurídico e temas irmãos) acha o tema certo, ou um parente
      direto, em 96\% dos casos contra 67\%. Esse ganho quase não apareceu nas questões
      publicadas (uma a mais), então não mudei a entrega.
\end{itemize}

\bd{O que recomendo mudar nos temas}
\begin{itemize}
\item Dar nomes que distingam os gêmeos (por exemplo, citar a disciplina no título). São oito
      pares, cujos números são os IDs na base de temas: 4684/6716, 5063/5466, 4688/6717,
      4702/6723, 4704/6724, 4703/6734, 4705/6725 e 4708/6728.
\item Fundir ou hierarquizar melhor os subtemas que repetem o conteúdo do pai.
\item Marcar os temas ``guarda-chuva'' (onde aceitar o nível geral é seguro) e os que exigem o
      nível específico.
\end{itemize}

\bd{O que preciso de você}
\begin{itemize}
\item Confirmar qual dos dois nomes é o oficial em cada um dos oito pares.
\item Um parecer sobre seis questões em que o modelo e o rótulo de IA discordam (primeiro o tema
      que o modelo publicou, depois o que o rótulo sugere):
      \begin{itemize}
      \item Q0120 (provas; resposta à acusação): Teoria Geral das Provas $\times$ Defesa Preliminar.
      \item Q0238 (reforço da fiança, CPP): Liberdade Provisória $\times$ Prisão em Flagrante.
      \item Q0273 (crime punido na modalidade culposa): Crime Doloso e Culposo $\times$
            Condescendência Criminosa.
      \item Q0016 (jurisprudência sobre prisão): Teoria Geral das Prisões $\times$ Prisão Preventiva.
      \item Q0306 (associação criminosa): Crimes Contra a Paz Pública $\times$ Organizações Criminosas.
      \item Q0085 (concussão, agente funcionário público): Concurso de Pessoas $\times$ Concussão e
            Excesso de Exação.
      \end{itemize}
\item Se puder, rotular de 20 a 30 questões em que as duas avaliações de IA discordam, para
      medirmos a precisão de verdade.
\end{itemize}

\end{document}
